% SPDX-FileCopyrightText: 2026 Yiyu Zhou <yiyu@yiyuzhou.io>
%
% SPDX-License-Identifier: GPL-3.0-or-later

\chapter{Outputs}
\label{ch:outputs}

A theater plays sound in more than one place, such as the house speakers, a
subwoofer, and the monitors on stage, which are often channels of one mixing
desk.  In cinq, each place is an output: a name that the show gives to one
channel, or two, of an audio device.  Cues send their sound to outputs, never
to devices or channels, so moving an output to other channels, or to another
device at another venue, moves every cue that sends to it.

A new show has one output, Main, on the first two channels of the default
device, and new cues send to it.  This chapter adds an output for each
speaker of your computer, then sends the rain of your own show to one and
then the other.

\section{Adding outputs}

Press \key{Control+o} and open \texttt{my-show.cinq.toml}.

\begin{enumerate}
\item Click \ui{Outputs} below the cue lists in the sidebar.  The outputs of
  the show take the place of the cue list, and the inspector shows the
  selected output, Main.
\item Press \ui{Add output}.  A new output appears below Main, on the first
  two channels of the default device.  In the inspector, set its \ui{Name} to
  \texttt{Left} and its \ui{Channels} to \texttt{1}.
\item Press \ui{Add output} again, and make an output named \texttt{Right} on
  channel~2.  The outputs now look as in Figure~\ref{fig:outputs}.
\end{enumerate}

\begin{figure}[H]
  \centering
  \includegraphics[width=\textwidth]{screenshots/outputs.png}
  \caption{Outputs of a show, which take the place of the cue list}
  \label{fig:outputs}
\end{figure}

Channels count from 1: \texttt{3} names one channel, and \texttt{1-2} names
two, left first.  The \ui{Device} list of an output chooses the device that
it plays on, such as a sound card or a desk with more channels.  Its default,
\ui{Default device}, is the device that the \ui{Default device} list in the
sidebar chooses for the whole computer, so a show made on one computer plays
on another.  An output that cannot play, such as one whose device is
unplugged, shows why in its \texttt{State} column: in yellow while cinq
tries the device again by itself, and in red where the output must change.
The outputs of one device share a warning in the sidebar, which shows them
when clicked, and cues play on the other outputs.

While the outputs show, the keys of cues act on them: \key{ArrowDown} and
\key{ArrowUp} select an output, and the keys that add, move, and delete cues
add, move, and delete outputs.

\section{Sending a cue}

\begin{enumerate}
\item Click \ui{Main} under \ui{Cue lists} in the sidebar to return to the
  cue list.
\item Select the rain, and click the \ui{Levels} tab.  If the exercise of
  Chapter~\ref{ch:fades} left its level at $-60$~dB, set it to $-20$~dB so
  that it can be heard.  Below the level, a field for each output holds the
  level at which the rain sends there: 0~dB for Main, and \texttt{off} for
  the others.
\item Type \texttt{0} in the field of Left, and empty the field of Main, as
  in Figure~\ref{fig:sends}.
\item Press \ui{Standby selected}, then \ui{GO}\@.  The rain plays from the
  left speaker alone.
\item While it plays, type \texttt{0} in the field of Right, and empty the
  field of Left.  The rain moves to the right speaker.
\end{enumerate}

\begin{figure}[H]
  \centering
  \includegraphics[width=\textwidth]{screenshots/sends.png}
  \caption{\ui{Levels} tab of a cue that sends to one output alone}
  \label{fig:sends}
\end{figure}

Each send applies after the level of the cue, which fade cues move, so a
fade moves every send of its target at once.  \ui{New cues send}, in the
inspector of an output, sets the level at which cues added later send
there.  Main has 0~dB and new outputs have none, so new cues send to Main
alone.  \ui{Edit outputs}, beside the fields, shows the outputs again.

\section*{Exercise}

You can now send each cue where it belongs.  While the rain plays, set the
\ui{Channels} of Right to \texttt{1}.  The rain moves to the left speaker,
although its cue has not changed.  Then put Right back on channel~2.
